\section{Data and methods}
\label{sec:methodology}

% ---------------------------------------------------------------------------
% Macros this revision's prose needs that analysis/emit_tex_values.py does not
% (red, visible) until the emit script reads the new results columns, and are
% silently superseded by the real \newcommand in values_generated.tex when it
% does. Nothing here defines a number. Declared in this section because it is
% the first of the four sections that use them.
% ---------------------------------------------------------------------------

The analysis proceeds in five steps. I first map macroeconomic oil and gas scenarios into a household price vector, shocking gas and electricity asymmetrically and pump prices separately. I then compute first-order household cost changes as quantity times price change on nationally grossed microdata --- on the consumption side, outside the tax-benefit calculator, for reasons set out below. Third, I cross the two incidence facts --- losses in pounds and losses as a share of income --- across the income distribution and across the finest geography the data will bear. Fourth, I score five competing policy responses on five reporting metrics and at \genEnvelopeRowKindCount{} row types. Fifth, I state what the design cannot recover. This section sets out each step, the microdata that carry it, and the limitations that bound the interpretation of every number in Section~\ref{sec:results}.

A methodological point is load-bearing enough to state at the outset. PolicyEngine UK, the microsimulation model used here, has no price channel: household energy quantities are imputed inputs, fixed by the data rather than determined by any price, and no policy parameter moves them. The shock is therefore applied on the consumption side rather than injected as a parameter reform, and three of the five policy responses are computed directly from the model's eligibility and household variables rather than as reforms. Subsection~\ref{sec:step2} documents exactly which objects are and are not present in the model. This is not a workaround reported reluctantly: the boundary of what an open tax-benefit model can express about an energy price shock is itself a result, it is precisely specified here, and it is the basis of the upstream contributions listed in Appendix~\ref{app:upstream}.

\subsection{Step 1: from a macro scenario to a household price vector}
\label{sec:step1}

Household energy prices in Great Britain are not market prices. For the roughly four in five households on a default tariff, the price per kilowatt hour is a regulated ceiling set quarterly by Ofgem, and that ceiling is a mechanical function of forward wholesale prices observed in an earlier averaging window, plus network, policy, operating and margin allowances. Any macro-to-micro exercise for the UK must therefore pass through the cap formula rather than through a spot price. I treat the cap as the transmission device, calibrate it outside the model, and carry the resulting retail price factors into Step 2 as a fixed vector of proportional price changes by fuel.

\paragraph{Scenarios.} I take three oil and gas paths. The \emph{baseline} follows the NIESR Summer 2026 Outlook \citep{niesr2026summer}, with Brent averaging \$103 per barrel and NBP gas consistent with that path; the \emph{adverse} case follows NIESR's alternative \citep{niesr2026spring} in which Brent reaches \$140 per barrel and gas and oil are sustained roughly 50 per cent above baseline; and the \emph{realised} case uses observed 2026 outturns to date --- UK gas peaking \genGasPeakPence{} pence per therm above its pre-war level and Brent peaking \genOilPeakUsd{} dollars per barrel above it, a \genOilPeakPct{} per cent move --- extended to the end of the modelling horizon at forward-curve levels. Reporting all three is deliberate: the realised path is the policy-relevant object for the Autumn Budget, while the baseline and adverse paths bound it and preserve comparability with the published macro forecasts that motivate the exercise. The realised path is the central case throughout.

\paragraph{Wholesale to cap.} The cap's wholesale allowance is built from forward contracts observed over an averaging window that closes shortly before each quarterly announcement and covers delivery in the subsequent charge period. Mechanically this places a lag of roughly four to five months --- one and a half quarters --- between the midpoint of a move in forward wholesale prices and the midpoint of the household bill that reflects it. The wholesale allowance itself is not the whole bill: on recent Ofgem decompositions it accounts for about 40 to 50 per cent of a typical dual-fuel bill, the remainder being network charges, environmental and social obligation costs, operating costs, debt-related allowances and the earnings-before-interest-and-tax margin, none of which respond to a commodity shock at the horizon considered here. The arithmetic I impose is therefore deliberately transparent. A proportional wholesale move of $+X$ per cent enters the cap as
\begin{equation}
\label{eq:passthrough}
\Delta \ln p^{\text{cap}}_{f,t} \;=\; \phi_f \, \Delta \ln p^{\text{wholesale}}_{f,t-\ell},
\end{equation}
with a fuel-specific pass-through coefficient $\phi_f$ and a lag $\ell$ of one and a half quarters. Two distinct attenuations are applied, and they are applied \emph{independently} rather than composed into a single reduced form. The first is the bill-level pass-through, $\phi_{\text{gas}} = \genBillLevelPassThrough{}$, which is the share of a proportional wholesale gas move that reaches a gas bill through the cap. The second is a gas-to-electricity ratio, $\psi = \genElecToGasPassThroughRatio{}$, applied to the electricity leg to reflect that gas sets the electricity price most but not all of the time and that the electricity cap embeds a larger share of non-commodity cost. An earlier statement of this method wrote $\phi_{\text{elec}} = \psi\,\phi_{\text{gas}}$ as though the two composed multiplicatively into a single coefficient; the implementation applies each attenuation to its own leg, and equation~\eqref{eq:passthrough} above is the corrected statement of what the code does. Neither coefficient is 0.45: that figure is the wholesale \emph{share} of a dual-fuel bill, which bounds $\phi$ from above but is not itself a pass-through. The lag is derived rather than assumed. Ofgem's observation window closes about seven weeks before the charge period begins and covers roughly the preceding quarter, so the midpoint of the window sits four to five months ahead of the midpoint of the bill it prices --- about one and a half quarters. Three quarters, which treats the window's closing date as its midpoint, is retained in the sweep of Appendix~\ref{app:caplag} rather than as the central case. Because the cap is reset quarterly and not continuously, the path is quantised into quarterly steps with a phase-in profile. That profile is constructed by building, for each cap period, the \emph{actual forward-curve observation window} Ofgem's methodology assigns to it, and computing what share of that window falls after the shock date of 28 February 2026. It is not a linear ramp slid along a calendar until it fits. The distinction is what makes the one-and-a-half-quarter lag an institutional property of the cap methodology rather than a parameter fitted to the answer, and it is why the sweep in Appendix~\ref{app:caplag} measures a residual rather than reproducing an identity: across lags of one to four quarters the mean loss moves by \pounds \genCapLagRangeGbp{} per household against a central \pounds \genCentralLossMean{}, or \genCapLagSpreadPct{} per cent. The annual price faced by a household in the microdata is the consumption-weighted average of the quarterly cap levels prevailing over the modelled window, not the peak and not the level the cap eventually reaches; the alternative, charging the steady-state level for the full twelve months, is retained and reported as the labelled \emph{steady-state} variant. It matters for magnitudes: the aggregate falls from \pounds \genSteadyStateAggBn{} billion to \pounds \genAggSpendRealisedBn{} billion when the phase-in is applied.

\paragraph{The modelled window.} Every annual figure in this paper is an average or a total over one named window: \textbf{March 2026 to February 2027}, the twelve months from the shock. The window is a single constant in the code, and the cap phase-in weights, the seasonal consumption weights and the pump damping fraction are all derived from it rather than written down separately. Annualising the two legs over different windows and summing them is not admissible, because calendar 2026 begins two months before the shock and ends before the cap has fully moved, so a domestic leg averaged over one window and a pump leg over another do not add to anything. Three windows were available. Calendar 2026 reports twelve months of denominator against ten months of shock and truncates the cap leg exactly where the cap's institutional lag puts it. The 2026Q4--2027Q3 window begins seven months after the shock and misses the pump-price peak entirely. The twelve months from the shock has neither defect: it contains one full heating cycle, so the seasonal weights sum to one without renormalisation, and it contains both the observed pump peak and the two cap resets --- October 2026 and January 2027 --- that the shock actually reaches. Its one cost is that it is not calendar 2026, and calendar 2026 is the window of the Resolution Foundation's \pounds 11 billion \citep{rf2026outlook}. That window is therefore retained as a separate run so the external comparison can be made like for like, and Section~\ref{sec:results-aggregate} states the window on both sides of it.

\paragraph{Anchoring the domestic leg: the pre-war cap is solved, not observed.} The counterfactual this paper measures against is what households would have paid absent the war, and no published cap is that object. The war began on 28 February 2026, and the forward-curve observation window behind the July 2026 cap of \pounds 1{,}663 closes about seven weeks before the charge period begins --- so that window already contained the war. Using \pounds 1{,}663 as a pre-shock baseline while simultaneously assigning the same quarter full war pass-through, which is what any profile phased on the shock date does, subtracts part of the shock from itself: the domestic leg is understated at source, and since the pump leg is not, the motor-fuel share is inflated by exactly the amount the domestic leg loses. This is a defect of identification and not of calibration, and correcting it moves more numbers in this paper than any other single change.

The pre-war level and the gas leg's peak-to-sustained fraction are therefore \emph{solved jointly} from both published caps. July 2026 (\pounds 1{,}663) and October 2026 (\pounds 1{,}723) each equal the pre-war level scaled by however much of the war had reached that period's own observation window, so two published caps give two equations in the two unknowns. One adjustment is required first: Ofgem's temporary removal of VAT on electricity holds the October cap roughly \pounds 45 below where the cap formula would otherwise put it, and that relief has nothing to do with the war, so it is added back. The solution is a pre-war counterfactual cap of \pounds \genCapBaselineLevel{} and a gas sustained fraction of \genGasSustainedFraction{} per cent, and it \genCapFitReproducesBoth{} reproduce \emph{both} published caps, to a maximum error of \pounds \genCapFitMaxErrorGbp{}. That residual is not evidence for the solve. Two equations in two unknowns leave zero degrees of freedom, so the error is zero by construction and nothing is over-identified; it establishes only that the solver converged. The evidence about the domestic leg is the fragility sweep below, not this residual. On that path the October cap moves \genCapAnchorMovePct{} per cent and July already carries \genCapBaseMovePct{} per cent, giving a war-consistent October level of \pounds \genCapAnnualised{}. Both figures are on the reduced typical-consumption basis introduced on 1 July 2026 \citep{ofgem2026tdcv}, which cut electricity from 2{,}700 to 2{,}500 kWh and gas from 11{,}500 to 9{,}500 kWh; the same October cap is \pounds 1{,}723 on that basis and \pounds 1{,}935 on the old one, so every cap figure in this paper is stated on the 2026 basis and comparisons across the July boundary are not made.

\paragraph{What the solve is fragile to.} Exact identification is not the same as robust identification, and the sensitivity that matters most in this paper is not the cap lag or the elasticity but the shape of the assumed wholesale gas peak. The solve takes the published caps as data and the monthly profile of the gas peak as an assumption; the second is not observed at the frequency the cap formula consumes it. Sweeping that monthly profile moves the solved pre-war counterfactual cap from \pounds \genGasProfileCapMinGbp{} to \pounds \genGasProfileCapMaxGbp{}, the aggregate loss from \pounds \genGasProfileAggMinBn{} billion to \pounds \genGasProfileAggMaxBn{} billion, and the motor-fuel share from \genGasProfileFuelShareMinPct{} to \genGasProfileFuelShareMaxPct{} per cent. \genGasProfileNonIdentifiedCount{} of the \genGasProfileCellCount{} cells of the sweep return no solution at all: the assumed profile puts so little of the peak inside the July and October observation windows that no pre-war level and sustained fraction jointly reproduce both published caps. That is the paper's real fragility. It is reported here and swept in Appendix~\ref{app:domesticleg} rather than carried as a caveat, because a range that wide on the composition result is a finding about what these data can identify, not a footnote to a point estimate.

One reporting convention needs stating because it is easy to slip. The coefficient $\phi_f$ is a share of the \emph{bill}, so the quantities the calibration delivers are bill uplifts, not unit-rate uplifts. There are also two distinct bill factors, and the paper reports the second. The \emph{steady-state} factor is the level the cap eventually reaches: \genGasUnitRatePct{} per cent above baseline on gas and \genElecUnitRatePct{} per cent on electricity. The \emph{annual applied} factor is what households actually pay over the modelled window, which is the steady-state factor scaled by the phase-in weight the window and the lag together imply --- \genAnnualPhaseInGasPct{} per cent on gas and \genAnnualPhaseInElecPct{} per cent on electricity. The annual applied factor, not the steady-state one, is what every result in Section~\ref{sec:results} is computed on. Because a standing charge does not move with wholesale costs, the implied \emph{unit-rate} increase is larger than either bill increase --- on the standing-charge shares in the current cap, materially so. Every figure quoted in this paper is a bill figure, and the elasticity sweep of Appendix~\ref{app:elasticity} is fed the same bill factor.

\paragraph{Asymmetric treatment of gas and electricity.} Gas and electricity are shocked \emph{separately and by different amounts}. Two institutional facts motivate this. First, gas is 62 per cent of UK final household energy consumption \citep{desnz2025ecuk}, the highest share among the European members of the G7 --- Italy, the next highest, is at 53 per cent --- so the direct gas channel is unusually large. No source ranks household gas shares across the G7 as a whole, since the comparable Eurostat series excludes the United States, Canada and Japan \citep{eurostat2018}, and the claim is stated at the strength the evidence supports. Second, under British marginal-cost wholesale market design a gas-fired plant sets the system marginal price roughly 85 per cent of the time, so a gas shock propagates into the electricity cap as well --- but attenuated, because the electricity cap embeds a much larger share of non-commodity costs and because the marginal-pricing share is below one. I therefore attenuate the electricity leg by a separate ratio $\psi = \genElecToGasPassThroughRatio{}$, calibrated to the product of the marginal-pricing frequency and the wholesale share of the electricity cap, applied independently of the bill-level pass-through $\phi_{\text{gas}} = \genBillLevelPassThrough{}$ rather than composed with it.

It is worth being precise about how much work this choice does, because the temptation to oversell it is real and the sensitivity analysis does not support it. Sweeping the marginal-pricing share from 0.70 to the naive symmetric value of 1.00 moves the mean household loss by \genAsymMeanShiftPct{} per cent and leaves the decile-one to decile-ten percentage ratio essentially frozen (Appendix~\ref{app:asymmetry}). The paper's distributional conclusions do not turn on this parameter. What the parameter does move, materially, is \emph{composition}: the gas share of the domestic-energy loss falls from \genGasShareDomesticLow{} per cent to \genGasShareDomesticHigh{} per cent across that range. That is the reason to keep the asymmetry. A gas-only social tariff, an electricity-side levy rebate and the JRF block's choice of which unit rate to discount all land differently depending on which fuel carries the loss, so fuel-level policy analysis is only meaningful with the two fuels separated --- but the headline incidence result is not an artefact of the separation, and I do not claim it is. PolicyEngine UK imputes gas and electricity as distinct household quantities rather than only as a combined domestic-energy total, which is what makes the separation implementable at household level; collapsing the two onto that total would discard the fuel split for no gain in the headline number.

\paragraph{Pump prices.} Road fuel is treated in parallel but more simply, because retail motor fuel is not capped and pass-through is fast --- pump prices track crude within weeks rather than quarters, so the motor-fuel channel carries no cap lag. The mapping from crude to pump is not derived from a duty-and-VAT accounting identity. What is done is simpler: the proportional pump-price move for each scenario is set directly from observed outturns --- roughly $+20$ per cent for petrol and $+36$ per cent for diesel at the 2026 peak --- and the scenario grid of Appendix~\ref{app:grid} then fits a pass-through coefficient to those stipulated points so that intermediate crude moves can be interpolated. The economics behind the stipulated points is the familiar one --- duty is specific rather than ad valorem, so a crude move passes into the pump price at less than one for one, while the VAT layer amplifies what does pass --- but it enters as a rationale for the calibrated numbers, not as an identity the model evaluates.

One calibration choice here is consequential enough to state explicitly. The observed pump-price moves are \emph{peaks}, and an annual incidence figure requires an annual average price, not a peak. The domestic side is already damped for exactly this reason: the wholesale gas peak enters the cap path through a sustained fraction of \genGasSustainedFraction{} per cent, which is not written down but solved so that the modelled cap reproduces Ofgem's confirmed October step \citep{ofgem2026cap}. Applying the peak pump move undamped for twelve months while damping the gas peak is not a defensible asymmetry. The stated justification for treating road fuel differently --- that pump prices pass through in weeks and the cap in quarters --- is a claim about \emph{lag}, not about \emph{duration}: it licenses applying the move sooner, not applying the peak for a full year. The main specification therefore damps the pump path on the same logic, at a sustained fraction of \genPumpSustainedFraction{} per cent, derived from the observed monthly pump profile over the same modelled window as the domestic leg. The two fractions are close --- and the gas fraction is now the \emph{larger} of the two. That is worth pausing on, because it reverses a relationship earlier drafts of this paper built a headline on. The intuition that a lagged forward-curve average must smooth a spike far more than a spot-tracking pump price does is not wrong as intuition, but on the solved pre-war baseline it is not what the data deliver: the war's effect on the forward curve was persistent enough that the cap carries most of it into the modelled window. On the defective observed-cap baseline the solve returned a gas fraction of roughly a fifth, and that number was absorbing the baseline error rather than measuring persistence. Because the motor-fuel share of the loss depends almost entirely on the \emph{ratio} of the two fractions, the correction removes most of the asymmetry that produced a large motor-fuel majority in earlier drafts. The undamped-peak variant is retained and reported throughout as an explicit upper bound, labelled the \emph{peak-fuel} specification, in which the aggregate loss is \pounds \genPeakFuelAggBn{} billion and the mean household loss \pounds \genPeakFuelLossMean{} against \pounds \genAggSpendRealisedBn{} billion and \pounds \genCentralLossMean{} in the main specification. Every fuel-share and gradient statement in Section~\ref{sec:results} is reported across both.

This channel turns out to carry a large part of the result on either calibration (Section~\ref{sec:results}), which is not what the framing of the policy debate would lead one to expect.

\subsection{Step 2: first-order incidence, computed on the consumption side}
\label{sec:step2}

\paragraph{Why the shock is not a parameter reform.} The natural implementation of a price shock in a tax-benefit microsimulation is a parameter reform: move the regulated price, let the model propagate it. That implementation is not available in PolicyEngine UK, and establishing precisely why is a contribution of this paper rather than an aside. Four features of the model, each verified against the release used here, close off the parameter-reform route.

First, household gas and electricity quantities and household petrol and diesel spend are all exogenous inputs, imputed from the Living Costs and Food Survey and calibrated to administrative consumption data. They are not functions of any price, so no change to a policy parameter moves them.

Second, the regulated retail price cap enters the model at only one point, and that point is the Energy Price Guarantee subsidy machinery rather than a household price channel. Raising the cap parameter therefore changes the size of a subsidy that has since lapsed; it does not change any household's bill.

Third, the model's contributed-parameter set includes a lever named for household energy bills --- the one whose name promises exactly this exercise --- but nothing in the model reads it, so setting it has no effect on any household.

Fourth, the pump price of road fuel \emph{is} parameter-driven, but because household road-fuel spend is an input, raising the price parameter holds spend fixed and reduces the implied litres. The model is thus implicitly unit-elastic in road fuel, and raising the pump price mechanically \emph{lowers} fuel duty receipts --- the opposite of a first-order price shock.

Finally, the Warm Home Discount is absent from the model altogether, as neither a parameter nor a household variable.

The shock is therefore applied directly to consumption, which is where a first-order price shock belongs in any case. For household $h$,
\begin{equation}
\label{eq:firstorder}
\Delta C_h \;=\; \sum_{f} q_{hf} \, \Delta p_{f},
\end{equation}
where $q_{hf}$ is the baseline quantity of fuel $f$ --- gas, electricity, petrol and diesel, kept separate --- and $\Delta p_f$ the modelled proportional price change from Step 1. This is the first-order approximation to the compensating variation familiar from \citet{deaton1980}: to a first order in prices the money-metric welfare loss equals pre-shock quantities times the price change, with the substitution term entering only at second order. Equation~\eqref{eq:firstorder} is therefore an explicit \emph{upper bound} on the true welfare loss, and I state it as such throughout. Holding quantities fixed also holds litres fixed, so fuel duty --- a specific, per-litre duty --- is unchanged by construction and only the VAT-inclusive spend rises, which is the internally consistent treatment and the one the model's own parameter path cannot deliver.

\paragraph{What the tax-benefit system is used for.} The shock itself never enters modelled household net income: energy spend is not a tax-benefit object in this model. PolicyEngine UK supplies the \emph{baseline} against which the shock and the policy responses are scored --- household net income, the LCF-imputed and NEED-calibrated energy quantities, the equivalisation and income-decile construction, the relative poverty indicators, and the survey weights that gross the sample to the household population. That is a substantial and non-trivial role, and it is what makes the exercise a microsimulation result rather than an arithmetic exercise on a survey. What the model does not supply is a price channel, and the paper does not pretend otherwise.

\paragraph{No estimated demand system.} I do not attempt to close the second-order gap with a demand system. A full QUAIDS specification in the spirit of \citet{banks1997} would in principle deliver price elasticities, but they are not identified in these data: the energy quantities are imputed rather than observed with independent price variation, there is no household-level price variation to exploit under a national cap, and cross-sectional expenditure variation identifies Engel curves rather than price responses. I therefore make the zero-elasticity assumption explicit and primitive, report the result as a bound, and relegate demand response to Appendix~\ref{app:elasticity}, where the pipeline is re-run under a flat elasticity grid and under five named published specifications. Measured as a welfare loss rather than as a change in expenditure, that sweep turns out to be a \emph{small} source of uncertainty --- demand response shaves at most \genWelfareShavedHighEpsPct{} per cent off the loss at the most elastic point of the grid. The large uncertainties in this paper are the fuel imputation, the accounting basis and the pass-through calibration, and the appendix now says so.

\paragraph{Reporting conventions.} Losses are reported two ways throughout: in pounds per year, and as a percentage of income. The income concept is equivalised household net income after housing costs on the households-below-average-income scale, which normalises a childless couple to one and weights the first adult at 0.58, each additional adult and each child aged fourteen or over at 0.42, and each child under fourteen at 0.20. This is the concept the published distributional literature uses, and it is the concept every percentage in this paper divides by. Section~\ref{sec:results-decile} reports the unequivalised gradient once as a robustness line so that the choice is visible rather than silent.

Percentages of income are the \emph{aggregate ratio} --- weighted loss divided by weighted income within the group --- rather than the mean of household-level ratios, because equivalised net income is zero or negative for a non-trivial part of the distribution and a mean of ratios is dominated by that tail. Households with non-positive equivalised income are dropped from every percentage statistic, and the dropped weight is reported wherever it is material. The tail is larger on the equivalised concept than on the unequivalised one: \genDecileOneZeroIncomeShare{} per cent of decile one, and all \genCoverageExcludedHouseholdsM{} million households (\genCoverageExcludedSharePct{} per cent) that fall outside deciles one to ten, which the model marks with a \(-1\) decile sentinel precisely because their equivalised income is non-positive. Where a decile-level share of the aggregate loss is reported, it is normalised on the total across the deciles shown rather than on a population total that includes households the table does not display.

\subsection{Step 3: crossing the two incidence facts}
\label{sec:step3}

The distributional literature contains two facts that are usually reported separately and that point in opposite directions. Budget-share analyses find that poor households spend a larger share of income on energy and therefore lose more in percentage terms. \citet{fetzer2024}, matching 22.2 million Energy Performance Certificates to National Energy Efficiency Data-Framework meter records, find the opposite ordering in cash: more affluent areas occupy larger, higher-consumption dwellings and lose more in pounds, which is why the 2022--23 Energy Price Guarantee, a per-unit subsidy, transferred disproportionately to them.

Both statements are true, and the policy interest lies in their intersection. I therefore compute both metrics on the same households and cross them: by income decile, within income decile, and by geography. The paired decile profile --- mean pound loss and mean loss as a share of income, on the same households and the same axis --- is the paper's headline object, and the divergence between the two orderings is what makes the compensation problem in Section~\ref{sec:policy} hard, because a policy calibrated to one metric necessarily misses the population identified by the other. The within-decile spread is reported alongside, since \citet{cronin2019}, \citet{douenne2020} and \citet{sallee2019} establish that horizontal dispersion routinely exceeds the vertical gradient.

\paragraph{Geography: region, not constituency.} An earlier design for this paper crossed the two metrics at parliamentary-constituency level. That is not producible from this data release, and the fact is worth recording precisely because the promise is common in this literature. The enhanced FRS release used here carries no constituency weight matrix --- it is calibrated to 149 area targets, not the 650 or 1{,}512 that a seat-level reweighting would require --- and its local-authority field is a degenerate default taking a single identical value for every household in the file. Region (twelve categories) and country (four) are the only genuine geography in the data, and they are what the paper reports. No constituency statistic appears anywhere in this paper, and none has been imputed to stand in for one.

\subsection{Step 4: scoring the policy responses}
\label{sec:step4}

Five candidate responses are scored against the shocked baseline. Three of the five do not exist in PolicyEngine UK and cannot be expressed as parameter reforms, so they are computed directly and transparently from the model's own eligibility and household variables. This is stated at the level of the individual instrument rather than buried, because it determines what each scored number means:

\begin{enumerate}
\item \textbf{Means-tested social tariff (35 per cent discount).} A proportional discount on the \emph{shocked} domestic energy bill for households receiving any means-tested benefit. \emph{Not in the model.} There is no social-tariff instrument; the Energy Price Guarantee machinery that superficially resembles one is a flat cap-ratio scaler with no means test. The means-tested population is constructed here from the eight means-tested entitlements the model carries --- Universal Credit, Pension Credit, child and working tax credits, Housing Benefit, Income Support, income-based JSA and income-based ESA --- flagged at household level.
\item \textbf{Universal discounted block (JRF).} A 50 per cent discount on the first half of \emph{typical} consumption, where typical consumption is the weighted median domestic bill, plus a \pounds 60 per-child allowance. The discount applies to the \emph{level} of the block, which is what JRF propose \citep{jrf2026}; discounting only the shock component would be a much smaller and quite different instrument. The modelled cost is \pounds \genJrfGapModelledBn{} billion against JRF's own \pounds \genJrfGapSponsorBn{} billion. The two are comparable --- both are sized on Ofgem's typical-consumption basis, which is JRF's --- and the gap is the discount \emph{rate}, which JRF do not publish and this paper set at \genJrfModelledDiscountPct{} per cent (Section~\ref{sec:results-policy}). Defining the block on a fixed quantity rather than as a proportion of each household's own bill is essential: a proportional discount would mechanically pay most to the largest users. The per-child allowance uses the model's own child count; an earlier implementation used household size minus two, which paid nothing to a lone parent with one child and paid an allowance to a three-adult household. \emph{Not in the model}: the residual Energy Price Guarantee machinery scales a single consumption allowance and has no two-tier block structure.
\item \textbf{Warm Home Discount expansion (\pounds 150).} A flat payment to the means-tested population, reflecting the 2026--27 scheme year reform reported to have scrapped the property-cost test and made the payment automatic. I have found no published Government document setting out that reform, so it is modelled as reported rather than as legislated. \emph{Not in the model at all}: no Warm Home Discount parameter and no Warm Home Discount variable exist in PolicyEngine UK.
\item \textbf{VAT zero-rating on domestic fuel and power, 5 per cent to zero.} LCF domestic energy spend is VAT-inclusive, so the gain is the VAT component of the shocked bill, $b\,(1 - 1/1.05)$. Note the ceiling this implies: the maximum relief is \genVatCutMaxOffsetPct{} per cent of a domestic bill, against applied bill rises of \genAppliedGasUnitRatePct{} per cent on gas and \genAppliedElecUnitRatePct{} per cent on electricity --- a factor of two to three rather than an order of magnitude --- and it reaches no motor fuel at all. The model's delta-only treatment of VAT (Section~\ref{sec:limitations}) is innocuous here because only the change is required.
\item \textbf{Flat per-household rebate (IPPR).} A uniform \pounds 183 credit funded by clawing back regulated network windfalls \citep{ippr2026}, with the financing side scored separately rather than netted into the household result.
\end{enumerate}

Each policy returns a per-household gain in pounds per year, set against the per-household loss from equation~\eqref{eq:firstorder}. Five metrics are reported, and the mix matters as much as any single one of them. Two are the conventional between-decile measures: \emph{cost per pound of bottom-decile gain}, total exchequer cost divided by the total gain accruing to decile one; and the \emph{share of spend reaching deciles one to three}. The first is a dimensionless ratio bounded below by one --- pounds of cost per pound of decile-one gain --- and carries no currency symbol anywhere in this paper, since a ratio of pounds to pounds is not a quantity of money. One is the within-decile count that motivates the exercise: the \emph{share of losers left uncompensated}, the fraction of losing households whose compensation falls short of their loss, reported overall and decile by decile. Two are continuous, and they are new to this version. The \emph{share of the aggregate loss offset} is total compensation delivered to losing households as a proportion of their total loss. The \emph{mean and median residual loss} is what is left of the loss after compensation, in pounds. The reason for adding them is that the uncompensated count is knife-edge: it treats a household left a pound short exactly as it treats one left five hundred pounds short, and it is applied to a loss that is itself an upper bound. An instrument can therefore look catastrophic on the count while offsetting more of the loss in aggregate, and leaving households closer to whole, than any of its rivals --- which is precisely what happens in Section~\ref{sec:results-policy}.

The instruments are also scored at a \emph{common exchequer envelope} of \pounds \genEnvelopeBn{} billion as well as at each sponsor's own stated cost, because scoring a \pounds \genVatCutCostBn{} billion instrument against a \pounds \genRebateCostBn{} billion one and reporting which leaves fewer households short measures generosity rather than design. But the common envelope has to be spent on something an instrument can actually do. Scaling each instrument's parameter until the envelope is exhausted produces rows no policy could occupy --- a \genSocialTariffScaledParameterPct{} per cent discount on the domestic bill, a \pounds \genWhdScaledParameterGbp{} Warm Home Discount, a negative rate of VAT --- and a saturation ranking generated at such parameters follows mechanically from paying households more than they lost. No conclusion in this paper rests on those rows.

The scorecard therefore carries \emph{\genEnvelopeRowKindCount{}} row types, and every row records which it is. Four are set out below; the fifth is the envelope-capped row, \genEnvelopeSemanticsCapped{}, which is reported separately from the feasible maximum precisely so that an instrument's own ceiling is not mistaken for the envelope's.

\begin{enumerate}
\item \emph{At the sponsor's own stated cost.} The instrument is scored at the parameter its sponsor proposes --- a 35 per cent social-tariff discount, a \pounds 150 Warm Home Discount, a 50 per cent discount on the block, a \pounds 183 rebate, five VAT points --- and the exchequer cost is whatever that comes to. This is the row that answers ``what does the proposal on the table do?'', and it is the only row on which the sponsors' published costings can be checked.
\item \emph{At the feasible maximum.} The instrument's own parameter is raised as far as \emph{the instrument itself} allows --- a 100 per cent bill discount; for the Warm Home Discount, a flat payment of \pounds \genWhdCeilingBillRuleGbp{}, which is the \emph{mean shocked domestic bill} of the eligible population --- the bill-exhausting rule, which is the one the code applies --- and not, as an earlier draft said, the payment that exhausts their loss, which is the loss-exhausting rule at \pounds \genWhdCeilingLossRuleGbp{}, smaller by a factor of \genWhdCeilingRuleRatio{} and costing \pounds \genWhdCeilingLossRuleCostBn{} billion against \pounds \genWhdCeilingBillRuleCostBn{} billion; both are stored, the paper quotes the bill-exhausting one throughout, and the difference is the motor fuel the means-tested population is imputed; the removal of all five VAT points --- and the envelope is whatever that costs. The ceiling is the instrument's, not the envelope's, and the two must not be run together. Where the instrument's ceiling falls short of \pounds \genEnvelopeBn{} billion, the shortfall is itself the result: VAT zero-rating can absorb only \pounds \genVatCutCostBn{} billion, the social tariff \pounds \genSocialTariffAbsorbableBn{} billion and the Warm Home Discount expansion \pounds \genWhdAbsorbableBn{} billion. Where it \emph{exceeds} the envelope the same rule applies in the other direction: the instrument's own ceiling is reported at its own cost, not truncated to the envelope. For the JRF block that ceiling is a 100 per cent discount, and it costs \pounds \genJrfFeasibleMaxCostBn{} billion --- several times the envelope, not equal to it. Table~\ref{tab:envelope} carries that row separately from the envelope-capped \genJrfCappedParameterPct{} per cent row, so the two ceilings --- the instrument's and the envelope's --- are never conflated.
\item \emph{At the scaled parameter.} The instrument is scaled past its feasible maximum until the envelope is exhausted. They are labelled by what they are rather than by the name of the policy they have left behind: ``a proportional domestic-bill subsidy at \genSocialTariffScaledParameterPct{} per cent, means-tested population'' rather than ``the social tariff'', ``a flat payment of \pounds \genWhdScaledParameterGbp{} a year to the means-tested population'' rather than ``the Warm Home Discount'', and ``a proportional domestic-bill subsidy at \genVatCutScaledParameterPct{} per cent, universal'' rather than ``VAT zero-rating''. Naming them generically is not a presentational nicety. A row that pays a household more than it lost produces a saturation ranking mechanically, and attaching a real policy's name to it is how an instrument that cannot exist comes to be crowned.
\item \emph{At widened eligibility.} The instrument's generosity is held exactly at the sponsor's own stated parameter and the envelope buys \emph{reach} instead, extending eligibility beyond the modelled means-tested population until the envelope is spent. This row is feasible in the sense that no parameter exceeds its own bound, but it is not administrable as modelled: households outside the means-tested population are admitted in ascending order of equivalised after-housing-costs income, which assumes the state can rank every non-claimant household on that concept. It is therefore reported as a perfect-observability \emph{upper bound} on what an envelope spent on reach could buy as \emph{targeting}. It is not an upper bound on the offset: \genSocialTariffAdmissionRuleCount{} admission rules are scored (Section~\ref{sec:policy-eligibility}), and the default rule turns out to offset the \emph{least} aggregate loss of the four, so the finding does not rest on the observability assumption. One of the two eligibility rows requires a label rather than a name, and the label is not a quibble: the Warm Home Discount's widened-eligibility row reaches an eligible share of \genWhdEligibleSharePct{} per cent of households. A \pounds 150 payment to every household is not a widened Warm Home Discount; it is arithmetically the flat universal rebate, at a different parameter, and it is reported as such.
\end{enumerate}

Comparing the second row type against the fourth is what isolates the margin on which an envelope is best spent, holding the instrument fixed and moving only whether the money buys generosity or coverage. Section~\ref{sec:results-policy} and Section~\ref{sec:policy} report that comparison as the paper's policy finding, and it connects directly to the coverage limitation below: an instrument conditioned on benefit receipt is bounded by how many households the condition can see, which in this model is \genMeansTestedHouseholdsM{} million against a Universal Credit caseload of 7.2 million. Where a sponsor has published a costing --- JRF's roughly \pounds \genJrfStatedCostBn{} billion for the block \citep{jrf2026}, IPPR's rebate \citep{ippr2026} --- the simulated cost is compared against it, and here the two \emph{are} defined on the same base: the block runs on Ofgem's typical-consumption basis by default, which puts the half-block at \pounds \genJrfBlockOfgemBasisGbp{}, and that is JRF's own basis. The alternative basis, this model's weighted median domestic bill at \pounds \genJrfBlockModelledGbp{}, is a labelled variant and not the headline. The comparison is therefore made rather than declined, and what it shows is a difference of one parameter rather than of costing method: JRF publish the block size and the total but not the discount rate, so the \genJrfModelledDiscountPct{} per cent scored here is this paper's assumption and JRF's own total implies \genJrfImpliedDiscountPct{} per cent on the same block (Section~\ref{sec:results-policy}). The version scored here is accordingly the more generous instrument, and it is generous by an assumption of this paper's rather than by anything JRF wrote.

\subsection{Microdata}
\label{sec:data}

All household results run on PolicyEngine UK, an open-source static tax-benefit microsimulation model of the United Kingdom, on its enhanced Family Resources Survey 2023--24 base: 535{,}080 microdata records grossing to \genTotalHouseholdsM{} million households, evaluated at period 2026. Three household totals appear in this paper and they are not inconsistent. \genTotalHouseholdsM{} million is the full grossed population. Table~\ref{tab:decile} sums to \genCoveredHouseholdsM{} million because \genCoverageExcludedHouseholdsM{} million households (\genCoverageExcludedSharePct{} per cent) carry non-positive equivalised income, fall outside deciles one to ten, and are dropped from every percentage statistic. \genTotalHouseholdsMFromMean{} million is the aggregate loss divided by the mean loss, which rounds differently. Each table states which of the three it is on. A reader who knows the FRS will read those two numbers as an error, since the survey interviews on the order of twenty thousand households a year, so the construction is worth stating. The enhanced dataset is not a sample of 535{,}080 households. It is the FRS sample replicated across a grid of imputation draws --- each survey household appearing many times with different draws of the imputed wealth, consumption and top-income variables --- with the survey weight divided across its own replicates, so that the weighted total is the household population and the replication carries the imputation distribution rather than only its mean. The effective sample is the FRS sample; the record count is a property of the imputation design. Variables absent from the FRS are imputed by quantile regression forests from donor surveys: wealth from the Wealth and Assets Survey; consumption from the Living Costs and Food Survey across twelve COICOP categories, with petrol, diesel, gas and electricity carried separately; a VAT expenditure rate from the Effects of Taxes and Benefits dataset; top incomes from the Survey of Personal Incomes; and capital gains following Advani and Summers. Quantile regression forests are used rather than conditional-mean imputation so that the imputed distributions, not merely their means, are preserved --- which matters here, because the object of interest is dispersion in energy quantities within income deciles.

Two features of the construction are load-bearing. First, the LCF-imputed energy variables are calibrated against National Energy Efficiency Data-Framework 2023 kilowatt-hour distributions, but they are carried through the model --- and through this paper --- as annual spending in pounds, not as physical quantities. Repricing them therefore means scaling a baseline expenditure by a proportional price factor, which is exactly equation~\eqref{eq:firstorder} with $q_{hf}\Delta p_f$ read as baseline spend times the proportional price change. They are not physical quantities, and nothing in the calculation depends on the distinction, since the first-order object is the same either way. Second, survey weights are re-optimised by gradient descent against administrative targets --- 149 of them in this release, at national and regional level. That is what supports the regional aggregation in Step 3 and, equally, what rules out the seat-level aggregation it replaced. Monetary variables are uprated to the modelling year on the OBR's March 2026 Economic and Fiscal Outlook. Prices in the baseline correspond to the Ofgem Q2 2026 cap.

\subsection{Limitations}
\label{sec:limitations}

The design has real and specific limits, and I state them rather than qualifying them.

\paragraph{The model has no price channel.} This is the most consequential structural limitation and it is documented variable by variable in Section~\ref{sec:step2}. Its consequences run in both directions. Because energy quantities are inputs, no shock can propagate endogenously: there is no substitution, no induced change in fuel duty, no feedback from prices into any other household variable, and the incidence computed here is arithmetic on a fixed basket rather than a model response. Because three of the five scored instruments do not exist in the model, their costings are the paper's own construction from the model's eligibility variables and not a reproduction of the model's own policy logic; a reader who disagrees with the parameterisation of the block or the social tariff is disagreeing with this paper, not with PolicyEngine. The offsetting virtue is that every step is explicit and auditable, and the gaps are precisely enough located to be fixed upstream.

\paragraph{No constituency geography.} A 650-seat cross of the two incidence metrics would be the most useful geography for the decision this paper addresses, and it is not deliverable from this data release: there is no constituency weight matrix, and the local-authority field is a degenerate default identical for every household. Region is therefore the finest geography reported, and no seat-level statistic is approximated in its place. Nothing in the paper should be read as a seat-level result.

\paragraph{No heating-fuel indicator, hence no off-gas-grid split.} The baseline variable set contains no primary-heating-fuel or gas-grid-connection variable, so the on-grid/off-grid split that \citet{douenne2020} makes central for France cannot be computed here. Roughly four million UK properties are off the gas grid, heating oil and LPG sit wholly outside the price cap, and these households were the group most often omitted from 2022 support schemes --- so this is a substantive gap, not a formality. It is reported as an open question in Section~\ref{sec:discussion} rather than answered with an imputation the data cannot support.

\paragraph{Modelled means-tested coverage misses the administrative caseload by a factor of two and a half.} All \genMeansTestedVarCount{} means-tested variables \emph{resolve} in the model --- none is silently missing --- but resolving and contributing are different things, and saying only the first is misleading. Four of the eight (child tax credit, working tax credit, income support and income-based jobseeker's allowance) contribute \emph{zero} households: they resolve to a variable that is zero for every record, being legacy benefits essentially fully migrated to Universal Credit by the modelling year. The means test in this paper is therefore effectively a four-benefit test. The modelled means-tested population is \genMeansTestedShareHouseholds{} per cent of households --- \genMeansTestedHouseholdsM{} million, comprising 3.03 million on Universal Credit, 1.22 million on Pension Credit, 0.73 million on Housing Benefit and 0.34 million on income-based ESA. The Universal Credit figure alone should be 7.2 million households on the Department for Work and Pensions caseload to May 2026: the model captures about two in five of the largest means-tested benefit in the country, and its total for \emph{all} means-tested benefits sits well below the caseload of that single benefit. This is not a take-up nuance; it is a data-quality finding about the model, reported here in the same spirit as the missing price channel, and it belongs in the upstream list of Appendix~\ref{app:upstream}.

The consequence for this paper is direct and one-sided. The reach of the social tariff and of the Warm Home Discount expansion is materially understated, their costs are understated, and their uncompensated shares are correspondingly \emph{overstated}. Those two rows of the scorecard are upper bounds on the uncompensated count, not estimates of it. One statistic is deliberately not reported anywhere in this paper: the share of households losing more than five per cent of net income that sit outside the means-tested system. With a means-tested population of \genMeansTestedShareHouseholds{} per cent of households, that share cannot arithmetically fall below about \genLargeLoserCeilingPct{} per cent whatever the incidence of the shock, so a figure in the high nineties carries almost no information about this shock and reads as though it did. It is dropped rather than qualified. The qualitative conclusion --- that a means test conditioned on benefit receipt leaves a large share of losers untouched --- is considerably more robust than any of the percentages attached to it.

\paragraph{The gradient is a cross-concept statistic: ranked on one income concept, measured on another.} This is a limitation and not a detail, and an earlier statement of it in this paper was wrong in the direction that flattered the result. Percentages of income divide by equivalised net income \emph{after} housing costs. The decile to which each household is assigned comes from the model's \texttt{household\_income\_decile} variable, and that variable has now been identified rather than guessed: it ranks on equivalised \emph{before}-housing-costs income, \emph{person}-weighted. Agreement with that concept is \genDecileRankAgreementPct{} per cent, person-weighted; agreement with the equivalised after-housing-costs concept the burdens actually use is \genDecileRankAgreementEquivAhcPrecisePct{} per cent. The earlier claim that the ranking matched the burden denominator rested on an audit that weighted households rather than people and tested the wrong candidate, and it is false: on the model's own definition the ranking variable \genDecileRankMatchesDenominator{} match the burden denominator.

Every decile gradient in this paper is therefore a cross-concept statistic, and should be read as one: households are sorted by before-housing-costs equivalised income and their burdens measured against after-housing-costs equivalised income. The two concepts differ most exactly where the gradient is steepest --- among low-income renters, whose housing costs are a large share of income --- so this is not a benign source of noise. The direction is not signable a priori: attenuation from ranking noise pushes the measured gradient down, while systematic movement of high-housing-cost households into lower after-housing-costs positions than their before-housing-costs rank pushes it up. Rebuilding the ranking on the burden concept is the single most useful robustness exercise this paper does not contain. The same variable also carries a \(-1\) sentinel for households with non-positive income, which is what accounts for the \genCoverageExcludedSharePct{} per cent of households discussed above as falling outside deciles one to ten.

\paragraph{The energy imputation is wrong in both directions, and the two errors compound.} Two calibration defects in the imputed consumption variables bear directly on the paper's central claim, and they run opposite ways, so their effect on the \emph{composition} of the loss is larger than the effect of either alone.

The first is motor fuel, which is over-imputed and over-imputed most at the bottom. The consumption fusion places \pounds \genFuelSpendDecileOneRawGbp{} of annual motor-fuel spend on the average decile-one household and \pounds \genFuelSpendDecileTenRawGbp{} on the average decile-ten household --- an essentially flat profile. ONS Family Spending for FYE 2025 \citep{ons2026familyspending}, the survey the imputation draws on, gives a far steeper one: \pounds \genOnsBenchFuelDecileOneGbp{} at the bottom against \pounds \genOnsBenchFuelDecileTenGbp{} at the top. One convention matters when reading those benchmarks: ONS publishes weekly figures, and they are annualised here at $365.25/7 = 52.18$ weeks rather than at 52, which is why the decile-ten benchmark is \pounds \genOnsBenchFuelDecileTenGbp{} and not the \pounds 1{,}357 a naive fifty-two-week conversion gives. The mechanism is visible in the participation margin, where the microdata give the top and bottom deciles an identical share of households recording any motor-fuel spend, whereas the Department for Transport's National Travel Survey table on household car \emph{availability} finds 40 per cent of the lowest income quintile with no car against 14 per cent of the highest \citep{dft2025nts}. Two qualifications on that comparator: it measures car availability rather than car access, and it covers England only, while the model is UK-wide. Neither disturbs the direction, and the direction is what matters here --- the imputation is loading motor fuel onto poor households that do not own cars.

The second is domestic energy, which is \emph{under}-imputed. Mean modelled domestic energy spend is \pounds \genModelDomesticMeanGbp{} against ONS Family Spending's \pounds \genOnsBenchDomesticMeanGbp{} for gas and electricity \citep{ons2026familyspending}, and against DESNZ's \pounds 1{,}899 for a fixed-consumption 2025--26 dual-fuel bill \citep{desnz2026qep}, a shortfall of roughly a quarter --- and below Ofgem's typical-consumption cap level, which is not a credible place for a population mean to sit, since the cap is a typical-consumption construct rather than an average bill. Smaller households, prepayment customers and households using no gas are all already inside the ONS mean, so the usual definitional explanations do not close the gap.

The two defects compound on the one result that most depends on them. The motor-fuel base is about a quarter too high and the domestic base about a quarter too low, so the motor-fuel \emph{share} of the loss inherits both errors in the same direction. This is one of the reasons Section~\ref{sec:results-fuel} reports that share as a range rather than as a point estimate, and it is why an ONS-consistent specification that corrects both levels is carried through the paper alongside one that corrects only the decile shape of motor fuel. Correcting the shape alone leaves the aggregate essentially unchanged, at \pounds \genOnsFuelAggBn{} billion, but cuts the decile-one to decile-ten burden ratio from \genDecileRatioPct{} to \genOnsFuelDecileRatioPct{} to one; correcting both levels gives \pounds \genOnsLevelsAggBn{} billion and a ratio of \genOnsLevelsDecileRatioPct{}, and cuts the motor-fuel share to \genOnsLevelsMotorFuelShareOfLoss{} per cent. Both are reweightings of the same households onto published aggregates and are therefore illustrative rather than corrected imputations --- ONS deciles are unequivalised gross-income deciles and these are equivalised net-income deciles --- but the direction is unambiguous and the size is too large to leave in a footnote. Both are also reported upstream, since they are properties of the model rather than of this application (Appendix~\ref{app:upstream}).

\paragraph{One live policy fact interacts with the scored instruments.} The domestic anchor is not itself an open issue: the pre-war cap and the gas sustained fraction are solved jointly against both published 2026 caps and reproduce each exactly (Section~\ref{sec:step1}). What remains open is that Ofgem's temporary removal of VAT on electricity from 1 October 2026 holds the October cap materially lower than it would otherwise be and breaks comparability with earlier quarters. That interacts directly with the VAT zero-rating instrument scored in Step 4: part of the relief that instrument is credited with delivering may already sit in the baseline this paper measures against, in which case the scored VAT row overstates the additional relief available. I flag this rather than adjust for it, because disentangling the two would require a cap decomposition Ofgem has not yet published.

\paragraph{The motor-fuel profile of means-tested households is an imputation, and the central policy claim is conditional on it.} The eligibility-over-generosity result of Section~\ref{sec:policy-eligibility} depends on how much of the aggregate loss the means-tested population carries, and that in turn depends on their motor-fuel spending, which is imputed rather than observed. The imputation gives means-tested households \pounds \genMeansTestedFuelSpendRawGbp{} of annual motor-fuel spending against \pounds \genNonMeansTestedFuelSpendRawGbp{} for everyone else, a ratio of \genMeansTestedFuelRatioRaw{}, with \genZeroFuelShareMeansTestedRawPct{} per cent recording none at all against \genZeroFuelShareOverallRawPct{} per cent. Some of that gap is real --- benefit receipt and car ownership are genuinely negatively correlated --- and some of it is the same participation defect that gives deciles one and ten identical fuel-purchasing rates, running the opposite way. This paper cannot separate the two. Two alternative specifications bound it: a \emph{parity} specification assigning means-tested households the population's own fuel profile, and a \emph{participation-corrected} specification matching their zero-fuel rate to external car-availability evidence while leaving conditional spending alone. The means-tested share of the aggregate loss moves from \genMeansTestedShareOfLossPct{} per cent to \genMeansTestedShareOfLossParityPct{} and \genMeansTestedShareOfLossParticipationPct{} per cent respectively, and Section~\ref{sec:policy-eligibility} reports the policy finding across all three.

\paragraph{The regional profile is not benchmarked.} The regional results in Section~\ref{sec:results-region} are reported because region is the finest geography the data supports, but they are not validated against any external source, and their shape invites the same suspicion as the decile profile. A mean loss in the North East well above Wales and several times London's, in the lowest-income English region, is more plausibly the motor-fuel imputation defect reappearing at regional level than a genuine feature of regional energy exposure. I report the regional profile as a description and build nothing on it.

\paragraph{No general equilibrium.} The model is static and use-side only. \citet{kanzig2023} finds that poorer households lose from energy shocks mainly through income and employment rather than through budget shares, and that general-equilibrium and indirect effects account for roughly two-thirds of the total response. \citet{goulder2019} make the complementary point that use-side effects are regressive while source-side effects are progressive. Taken together, the incidence reported here is a component of total incidence, plausibly around a third of it, and the direction of the omitted component is not guaranteed to reinforce the reported one. Every headline figure should be read as a lower bound on total household incidence and an upper bound on the use-side channel conditional on zero substitution.

\paragraph{No consumption elasticity in the model.} PolicyEngine UK has no demand elasticity block and no pass-through module, a gap its maintainers have documented and which is noted against UKMOD's own practice of assuming an elasticity of 0.8. The main specification is zero-elasticity by construction and elasticity enters only through the external sweep in Appendix~\ref{app:elasticity}. Relatedly, CPI, CPIH and RPI exist in the model only as uprating indices, so the price path of Step 1 is necessarily constructed outside the model.

\paragraph{VAT enters only as a delta.} The model's household tax concept includes the \emph{change} in VAT under a reform but not the baseline VAT liability itself. Baseline VAT is never deducted from net income, so levels of net income are overstated relative to a full indirect-tax concept, and only reforms to VAT --- not the baseline burden --- affect the distribution. Fuel duty and the carbon tax module do enter at level, so the treatment across indirect taxes is inconsistent. For the VAT reform scored in Step 4 only the delta is required, so that reform result is unaffected; the baseline distribution is affected.

\paragraph{The VAT coverage factor.} The VAT imputation rests on a 0.38 coverage grossing factor, following IFS TAXBEN, applied on top of a four-predictor imputation from the Effects of Taxes and Benefits dataset. The factor is not independently validated, and is flagged as an outstanding issue by the model's maintainers.

\paragraph{Missing duties.} Alcohol, tobacco, air passenger duty, insurance premium tax and vehicle excise duty are absent from the model. This matters less for an energy shock than it would for a general indirect-tax study, but the reported indirect-tax base is incomplete and the compensation metrics do not net out indirect-tax changes outside energy.

\paragraph{Static, no lifecycle.} There is no intertemporal margin: households cannot borrow against a temporary shock, run down savings, or accumulate arrears. Given that energy arrears are the most visible real-world manifestation of the 2026 shock (Section~\ref{sec:discussion}), this is a substantive omission rather than a formality.

\paragraph{Data vintages.} Wealth is imputed from Wealth and Assets Survey round 7, covering 2018--20, which predates both the 2022 shock and the subsequent rate cycle; \citet{pallotti2024} show that nominal net asset positions rather than energy baskets drive most household heterogeneity in the euro-area response, so a stale wealth distribution is a first-order concern for any extension in that direction. Benefit take-up is modelled as a scalar rather than a behavioural model.

\paragraph{Access.} The enhanced microdata are not publicly redistributable and require authorised access, so the full pipeline cannot be re-run without it, although the code, parameters and aggregate results underlying every number reported here are public. Appendix~\ref{app:reproduce} sets out the arrangements.
